\scanpage{333}
\mspage{467}

where \deleted{\illegible{1}}
\[
(48)\qquad
\unsure{(\varepsilon_0)_0}=-\sigma_0^{-1}\rho_0=
\begin{pmatrix}1&\dfrac{1}{\unsure{v|\nu}}\\0&1\end{pmatrix},\qquad
\deleted{\unsure{(\varepsilon_1)_0}}
\inserted{(\varepsilon_1)_0=-\rho_0\sigma_0^{-1}}=
\begin{pmatrix}1&0\\1&1\end{pmatrix}
\]
--- and conversely, for $\rho_0,\sigma_0$ as above, and any isomorphism
$\varphi:G_0=\mathrm{Gl}(2)_S\xrightarrow{\sim}G$, the images
$\widetilde\rho,\widetilde\sigma$ of $\rho,\sigma$ in $PG$ satisfy
$(41)$, and the images $\rho,\sigma$ of $\rho_0,\sigma_0$ under $\varphi$
are liftings of $\widetilde\rho,\widetilde\sigma$ satisfying $(44)$.

\uline{Modular pinning theorem.} Let $G,SG,PG$ be forms
(\uuline{over} scheme $S$) of $\mathrm{Gl}(2),\mathrm{Sl}(2),\mathrm{GP}(1)$,
\unsure{associated} with an Azumaya \unsure{algebra} $M$ of rank $4$.

$1.)$ Let us call a \uline{modular \deleted{\illegible{1}}\inserted{triple}}
in $PG$ a triple of sections $(\widetilde\rho,\widetilde\sigma,
\widetilde\varepsilon_0)$ of $PG$, satisfying conditions $(49)(50)(51)$.

\deleted{$(\unsure{49})\quad \widetilde\rho^3=\widetilde\sigma^2=1$, $\widetilde\varepsilon_0$ is \unsure{regular} unipotent.}

\[
(49)\qquad \widetilde\varepsilon_0=\widetilde\sigma^{-1}\widetilde\rho
\quad\text{i.e.}\quad
\widetilde\rho^{-1}\widetilde\sigma\widetilde\varepsilon_0=1
\]
[so that
\[
(49\ \mathrm{bis})\qquad
\widetilde\varepsilon_1\overset{\mathrm{dfn}}=
\widetilde\rho\widetilde\sigma^{-1}=
\widetilde\rho(\widetilde\varepsilon_0)=
\widetilde\sigma(\widetilde\varepsilon_0)
\]
];

$(50)\qquad \widetilde\rho^3=\widetilde\sigma^2=1$,\quad
$\widetilde\varepsilon_0$ is regular unipotent
($\Longleftrightarrow\widetilde\varepsilon_1$ is \uline{regular} unipotent);

$(51)\qquad \widetilde\varepsilon_0$ and $\widetilde\varepsilon_1$
\deleted{\unsure{everywhere}} inequivalent \deleted{\unsure{on the}}
\inserted{in each} fiber
(which means, introducing the liftings
\[
\varepsilon_0=1+N_0,\qquad \varepsilon_1=1+N_1
\]
\margin{($\unsure{N_0\text{ and }N_1}$ strictly nilpotent,\ldots\ cf. $(42)$)}
to $G$, that $N_0$ and $N_1$ not proportional in every fiber).

$2)$ Let us call a modular triple \unsure{of|in} $G$ a
\mspage{332}

\scanpage{334}

triple of sections $(\rho,\sigma,\varepsilon_0)$ satisfying the
\deleted{\unsure{five}} following conditions

\[
\begin{aligned}
(52)\qquad &-\sigma^{-1}\rho=\varepsilon_0
\quad(\Longleftrightarrow\ \rho^{-1}\sigma\varepsilon_0=-1.)\\
&\left[\begin{aligned}
&\Longleftrightarrow\ -\rho\sigma^{-1}=\varepsilon_1
&&\text{where }\varepsilon_1=\rho\varepsilon_0\rho^{-1}\\
&\Longleftrightarrow\ -\rho\sigma^{-1}=\varepsilon_1
&&\text{where }\varepsilon_1=\sigma\varepsilon_0\sigma^{-1}
\end{aligned}\right]
\end{aligned}
\]
\[
(53)\qquad(\rho+\sigma)^2=\rho+\sigma
\]

$(54)\qquad\varepsilon_0$
\deleted{and $\varepsilon_1\overset{\mathrm{dfn}}=\rho\varepsilon_0\rho^{-1}=\sigma\varepsilon_0\sigma^{-1}$ \unsure{are}}
\inserted{is regular unipotent i.e. $\varepsilon_0=1+N_0$ with}
$N_0$ \deleted{\unsure{unipotent}} regular nilpotent
($\mathrm{Tr}\,N_0=\det N_0=0$, $N_0$ nonzero in each fiber).

\margin{$\Longleftrightarrow\left\{\begin{array}{l}
\det\varepsilon_0=1\\
\mathrm{Tr}\,\varepsilon_0=2\\
\varepsilon_0\ne1\text{ in every fiber}
\end{array}\right.$}

($\Longleftrightarrow\ \deleted{\unsure{\varepsilon_0}}\ \varepsilon_1
\overset{\mathrm{dfn}}=\rho\varepsilon_0\rho^{-1}
=\sigma\varepsilon_0\sigma^{-1}$ is regular unipotent)

$(55)\qquad\deleted{\rho}\cdot\varepsilon_0$ and $\varepsilon_1$ are
inequivalent in each fiber
($\Longleftrightarrow\ N_0$ and $N_1$ not proportional in every fiber).

$(56)\qquad$ We have one of the four equivalent conditions
\[
\begin{array}{cc}
\boxed{\sigma^2+1=0}&\boxed{\rho^2\unsure{+|-}\rho+1=0}\\[6pt]
\boxed{\mathrm{Tr}\,\sigma=0,\det\sigma=1}&
\boxed{\mathrm{Tr}\,\rho=\det\rho=1}
\end{array}
\quad.
\]

\deleted{\unsure{3}} $3)$ Let us call the \uline{\unsure{standard}}
\uline{modular triple} in $\mathrm{Gl}(2)_S$
(resp. in $\mathrm{GP}(1)_S$) the triple
\[
(57)\quad\left(\rho_0=\begin{pmatrix}0&1\\-1&1\end{pmatrix},\quad
\sigma_0=\begin{pmatrix}0&-1\\1&0\end{pmatrix},\quad
\varepsilon_0=\begin{pmatrix}1&-1\\0&1\end{pmatrix}\right)
\]
(resp. its image $(\widetilde\rho_0,\widetilde\sigma_0,
\unsure{\widetilde\varepsilon_0})$ in $\Gamma\mathrm{GP}(1)_S$) ---
which is indeed a modular triple.

This being so, we have \deleted{some}\inserted{the \unsure{commutative} diagram of} canonical
maps
\mspage{333}

\scanpage{335}
\mspage{468}

\[
\varphi\longmapsto(\varphi(\rho_0),\varphi(\sigma_0),\varphi(\varepsilon_0))
\]

$(58)$
\diagram{\begin{tikzcd}[ampersand replacement=\&,column sep=large,row sep=large]
\operatorname{Isom}(\mathrm{Gl}(2),G)
\arrow[r]\arrow[d,phantom,"\unsure{\wr}"]\&
\mathrm{TripMod}(G)\arrow[d]\&
(\rho,\sigma,\varepsilon_0)\arrow[d,mapsto]\\
\operatorname{Isom}(\mathrm{GP}(1),PG)\arrow[r]\&
\mathrm{Tripmod}(PG)\&
(\widetilde\rho,\widetilde\sigma,\widetilde\varepsilon_0)
\end{tikzcd}}
\[
\psi\longmapsto(\psi(\widetilde\rho_0),\psi(\widetilde\sigma_0),
\psi(\widetilde\varepsilon_0))
\]

and we have the

\uline{Pinning theorem.} All the arrows of $(58)$ are bijections.

This theorem only summarizes \deleted{the}\inserted{some}
\deleted{\unsure{reflections}}\inserted{of the} main results of the reflections
\unsure{of the|of the} present \S\ \uline{$\overline{\mathrm{VI}}$}.
Let us note that the injectivity of the horizontal arrows means \inserted{also}
that an \deleted{\unsure{isomorphism}}
$\mathrm{Gl}(2)_S=G_0\longrightarrow G$
(resp. $\mathrm{GP}(1)_S=PG_0\longrightarrow\unsure{G|PG}$
is known when we know its effect on $\sigma$ and $\rho$
(thus also on $\varepsilon_0$]
\unsure{$\leftarrow$} resp. when we know its effect on
$\widetilde\sigma$ and $\widetilde\rho$) --- which also means
the \unsure{same|=} thing for the automorphisms of
$\mathrm{Gl}(2)_S$ (resp. of $\mathrm{GP}(1)_S$)
which also means that
\[
(59)\qquad T_\rho\cap T_\sigma=\deleted{\unsure{\mu_2|\{1\}}}\ G_m
\]
(the intersection of the centralizers of $\rho,\sigma$ in
$\mathrm{Gl}(2)_S$ is \unsure{reduced|reduced} to the center $G_m$) ---
\unsure{resp.} the stronger relation
\[
(60)\qquad N'_\rho\cap N'_\sigma=G_m,
\]
\mspage{334}

\scanpage{336}

where $N'_\rho$, $N'_\sigma$ are the group \unsure{subschemes} of $\mathrm{Gl}(2)_S$ defined \deleted{by} \inserted{\unsure{respective\textsuperscript{ly}}} as stabilizers of $\widetilde\rho$, $\widetilde\sigma$ in $\unsure{\mathrm{Gl}(2)_S}$, i.e.
\[
(61)\qquad
\left\{\begin{aligned}
N'_\rho&=\{g\mid g(\rho)=\varepsilon\rho,\ \varepsilon\text{ scalar}\},\\
N'_\sigma&=\{g\mid g(\sigma)=\varepsilon\sigma,\ \varepsilon\text{ scalar}\}
\end{aligned}\right.
\]
\margin{N.B. we \unsure{will have} \unsure{necessari\textsuperscript{ly}} $\varepsilon^2=1$}

\uline{Remark} It is plausible, \deleted{that \illegible{5}} \inserted{\unsure{In view of these results,}} that $\mathrm{GP}(1)_{\mathbb Z}$ --- as an ``absolute'' sheaf of groups (over $\operatorname{Spec}\mathbb Z$) for fppf, is generated by the homomorphisms
\[
\begin{aligned}
\mathbb Z/2\mathbb Z&\xrightarrow{\alpha\,\mapsto\,\widetilde\sigma_0^{\alpha}}\mathrm{GP}(1)_{\mathbb Z}\\
\mathbb Z/3\mathbb Z&\xrightarrow{\beta\,\mapsto\,\widetilde\rho_0^{\beta}}\mathrm{GP}(1)_{\mathbb Z}\\
\mathbb G_a&\xrightarrow{\lambda\,\mapsto\,\widetilde\varepsilon_0(\lambda)}\mathrm{GP}(1)_{\mathbb Z}
\end{aligned}
\]
satisfying \unsure{the|the} \inserted{sole} relation (\unsure{R})
\[
(62)\qquad \widetilde\rho_0^{-1}\widetilde\sigma_0\widetilde\varepsilon_0(1)=1
\qquad\deleted{\unsure{(?)}}
\]
--- it would be well worth the trouble to clear up this question, as well as the analogous question for $\mathrm{Gl}(2)_{\mathbb Z}$, where the \deleted{relations} groups $\mathbb Z/2\mathbb Z$, $\mathbb Z/3\mathbb Z$ must \unsure{be replaced} by $\mathbb Z/4\mathbb Z$, $\mathbb Z/6\mathbb Z$, and the relation (62), by
\[
(63)\qquad \rho_0^{-1}\sigma_0\widetilde\varepsilon_0(1)=\rho_0^3=\sigma_0^2.
\]
\mspage{335}

\scanpage{337}
\mspage{469}

\uline{Remark} I have the distinct impression that the package of conditions (51) to (56) is redundant, and wonder for example whether \unsure{the set|one} of the three simple algebraic conditions
\[
(64)\qquad \sigma^2+1=\rho^2-\rho+1=0,\qquad (\rho+\sigma)^2=\rho+\sigma
\]
does not \unsure{imply} the other conditions \inserted{(defining $\varepsilon_0$ by (\unsure{51|52}))}. The third condition (64) means that
\[
(65)\qquad N=\rho+\sigma
\]
is \unsure{idempotent}
\[
(66)\qquad N^2=N,
\]
and in terms of $N$ we will have
\[
\sigma^{-1}N=\sigma^{-1}\rho+1\quad\text{i.e.}
\]
\[
(67)\qquad
\left\{\begin{aligned}
\varepsilon_0&\overset{\mathrm{dfn}}{=}-\sigma^{-1}\rho=1-\sigma^{-1}N=1+N_0
&\qquad N_0&=-\sigma^{-1}N\\[-2pt]
&&&=\sigma N\\
\varepsilon_1&\overset{\mathrm{dfn}}{=}-\rho\sigma^{-1}=1-N\sigma^{-1}=1+N_1
& N_1&=-N\sigma^{-1}=N\sigma
\end{aligned}\right.
\]
\margin{\uline{N.B.} We must \deleted{\unsure{perhaps}} \uline{specify} \deleted{\unsure{that in}} \illegible{4} \unsure{first} relations \unsure{(64)} \illegible{3}
$\operatorname{Tr}\sigma=0$, $\det\sigma=1$, $\operatorname{Tr}\rho=\unsure{0|1}$, $\det\rho=1$ \unsure{must require that} \illegible{2} are \uline{regular} i.e. \unsure{not homotheties} in any fiber \illegible{1}.}

and the condition $\varepsilon_0$ \deleted{\unsure{strict\textsuperscript{ly} unipotent}} \inserted{\unsure{regular unipotent}} is expressed by the conditions
\[
(68)\qquad \det\sigma N=\operatorname{Tr}\sigma N=0,\quad
\begin{gathered}
\sigma N\text{ nonzero in every fiber.}\\
\text{i.e. }N\text{ nonzero in every fiber}
\end{gathered}
\]
Now the \deleted{conditions} \inserted{third of these} (68) follows from the others, for if we had $N=0$ on a fiber, we could suppose $S$ the spectrum of a field, $N=0$, i.e. $\rho=-\sigma$, so \deleted{\illegible{1}} we would have both
\[
\sigma^2+1=0\qquad \sigma^2+\sigma+1=0
\qquad\text{whence }\sigma=0,\text{ absurd.}
\]
I also say that in no fiber do we have $N=1$, otherwise we would have in this fiber \deleted{$\rho=1-\sigma$} $\sigma=1-\rho$ and substituting in $\sigma^2+1=0$ we find \deleted{\illegible{1}}
\[
(1-\rho)^2+1=0\quad\text{i.e.}\quad \rho^2-2\rho+2=0,
\]
\mspage{336}

\scanpage{338}

comparing with $\rho^2-\rho+1=0$, we find \deleted{$\unsure{\rho-1}$} by subtraction $\rho-1=0$ i.e. $\rho=1$, but the relation $\rho^2-\rho+1=0$ becomes $1=0$, absurd.

Thus the idempotent operator $N$ is \deleted{everywhere} of rank $1$ \inserted{in each fiber}, so \illegible{1} (unless I am mistaken) \deleted{\unsure{says}} that $N$ is \deleted{decomposable} (in terms of $\unsure{\mathcal M}\simeq\check F\otimes F$) \inserted{local\textsuperscript{ly}} decomposable in the form $x'\otimes x$, with $\langle x,x'\rangle=1$, and \deleted{$\operatorname{Tr}$}
\[
\operatorname{Tr}N=1,\quad\det N=0
\]
so the relation $\det(\sigma N)=0$ in (68) is also automatic, since $\det\sigma N=\det\sigma\det N=0$. There remains only the relation
\[
\operatorname{Tr}\underbrace{\sigma N}_{\substack{\sigma^2+\sigma\rho\\=\sigma\rho-1}}=0
\]

( \deleted{that} which is also written
\[
(69)\qquad \operatorname{Tr}\sigma\rho=2\quad(\Longleftrightarrow\operatorname{Tr}\sigma^{-1}\rho=-2)
\]
\unsure{for} which we would have to see whether or not it follows from the other relations.

Let us take
\[
(70)\qquad
\left\{\begin{aligned}
\rho&=\rho_0=\begin{pmatrix}0&1\\-1&1\end{pmatrix}\\
\sigma&=\underline u\sigma_0\underline u^{-1}
=\begin{pmatrix}\unsure{a}&b\\c&d\end{pmatrix}
\begin{pmatrix}0&-1\\1&0\end{pmatrix}
\begin{pmatrix}d&-b\\-c&a\end{pmatrix}\\
&\text{where }\underline u=\begin{pmatrix}a&b\\c&d\end{pmatrix}\in\unsure{\Gamma\mathrm{Sl}(2)},\text{ i.e. }ad-bc=1
\end{aligned}\right.
\]
We find
\mspage{337}

\scanpage{339}
\mspage{470}

\[
(70\mathrm{bis})\qquad
\sigma=\begin{pmatrix}
bd+ac&-(a^2+b^2)\\
c^2+d^2&-(bd+ac)
\end{pmatrix}
\]
whence
\[
N=\rho+\sigma=\begin{pmatrix}
(bd+ac)&1-(a^2+b^2)\\
-1+(c^2+d^2)&1-(bd+ac)
\end{pmatrix}
\]
\[
N-1=\begin{pmatrix}
(bd+ac)-1&1-(a^2+b^2)\\
-1+(c^2+d^2)&-(bd+ac)
\end{pmatrix}
\quad\begin{gathered}
=-\det(N)N^{-1}\\
\text{if }\det N\text{ \unsure{invertible}}
\end{gathered}
\]
and we find
\[
(71)\qquad N^2-N=N(N-1)=-\det(N).1
\]
so \deleted{the condition}
\[
(72)\qquad N^2=N\deleted{\unsure{\Longleftrightarrow}}\quad\Longleftrightarrow\quad\det N=0\ ;\quad(\Longrightarrow\operatorname{Tr}N=1)
\]
Now we have
\[
\begin{aligned}
\det N={}&bd+ac-(bd\unsure{+}ac)^2\deleted{\unsure{+1}}\\
&\unsure{-|=}\bigl[-1-(a^2+b^2)(c^2+d^2)+a^2+b^2+c^2+d^2\bigr]\\
={}&bd+ac-\deleted{b^2d^2}-\deleted{a^2c^2}-2a\uline{bc}d+1
+\deleted{a^2c^2}+\uline{a^2d^2}+\uline{b^2c^2}+\deleted{\unsure{b^2d^2}}\\
&\hspace{4em}-a^2-b^2-c^2-d^2\\
={}&(bd+ac)-(a^2+b^2+c^2+d^2)+\underset{1}{\underbrace{(ad-bc)^2}}+1\\
={}&2+(bd+ac)-(a^2+b^2+c^2+d^2)
\end{aligned}
\]
so the condition \unsure{(72)|(71)} is equivalent to
\[
(73)\qquad a^2+b^2+c^2+d^2=2+(bd+ac),
\]
on the other hand we have by 70, (70bis)
\[
\unsure{\rho\sigma}=\begin{pmatrix}
c^2+d^2&-(bd+ac)\\
c^2+d^2-(bd+ac)&a^2+b^2-(bd+ac)
\end{pmatrix}
\]
so
\[
(74)\qquad \operatorname{Tr}\unsure{\sigma\rho}=\operatorname{Tr}\rho\sigma
=a^2+b^2+c^2+d^2-(bd+ac)=2
\]
\mspage{338}